%! AP Statistics — Unit 8, Lesson 8.3 — Warm-Up (Answer Key)
\documentclass[10pt]{article}
\usepackage{apstats-article}
\usepackage{apstats-key}

\begin{document}

\pageheader{Unit 8, Lesson 8.3}{Warm-Up --- Answer Key}
\namedateperiod

\vspace{0.15in}

A school cafeteria offers 4 lunch options. Administrators claim each option is equally
popular among students. A random sample of 80 students is asked to choose their
preferred lunch option. The results are shown below.

\renewcommand{\arraystretch}{1.5}
\begin{center}
\begin{tabular}{lcccc}
    \textbf{Option}   & A  & B  & C  & D  \\
    \hline
    \textbf{Observed} & 25 & 18 & 22 & 15 \\
\end{tabular}
\end{center}

\vspace{0.1in}

\begin{enumerate}

  \item State the null and alternative hypotheses for a chi-square goodness-of-fit test.

        \vspace{0.1in}
        $H_0$: \ans{Each lunch option is chosen by $1/4$ of all students
        ($p_A = p_B = p_C = p_D = 0.25$).}

        \vspace{0.05in}
        $H_a$: \ans{At least one option is chosen by a different proportion of students.}

        \vspace{0.1in}

  \item Calculate the expected count for each lunch option. Show your work.

        \vspace{0.08in}
        $E = n \times p_0 =$ \ans{$80$} $\times$ \ans{$0.25$} $=$ \ans{$20$}

        \vspace{0.1in}
        Each option has an expected count of \ans{20} students.

        \vspace{0.1in}

  \item Verify that the large-counts condition is satisfied. State your conclusion in a
        complete sentence.

        \ansline{All expected counts equal 20, which is at least 5, so the large-counts
        condition is satisfied.}

\end{enumerate}

\end{document}
